\documentclass{article}
\usepackage[UTF8]{ctex}
\usepackage{amsmath,amssymb,amsfonts}
\usepackage{enumitem}
\usepackage{geometry}
\geometry{a4paper, margin=1in}
% === 错题本排版优化（不参与题目标识） ===
\setlength{\parindent}{0pt}
\setlength{\parskip}{0.35em plus 0.10em minus 0.08em}
\setlength{\emergencystretch}{3em}
\setlist{itemsep=0.25em, topsep=0.40em, parsep=0pt}
\allowdisplaybreaks[2]
% === 错题本排版优化结束 ===
\begin{document}
\section*{上海大学2026年数学分析真题}
\section*{47．上海大学2026年数学分析试卷（回忆版）}

\subsection{题目1}
1．求极限：

\begin{itemize}
  \item[（1）] \(\lim _{n \rightarrow \infty} \frac{7^{n} n!}{(3 n)^{n}}\)
  \item[（2）] \(\lim _{n \rightarrow \infty}\left(\frac{1}{3}+\frac{2}{3^{2}}+\cdots+\frac{n}{3^{n}}\right)\) ．
\end{itemize}

\subsection{题目2}
2 ．求积分 \(\int_{0}^{\frac{\pi}{2}} \frac{1}{\sin x} \ln \frac{1+a \sin x}{1-a \sin x} \mathrm{~d} x\) ．

\subsection{题目3}
3．求 \(x^{2}+y^{2}=2 a z\) 与 \(\left(x^{2}+y^{2}\right)^{2}=2 a^{2} x y\) 相交区域的体积．

\subsection{题目4}
4．设

\[
f(x, y)=(x-y)^{2}+\iint_{\sum} x\left(z^{2}+e^{z}\right) \mathrm{d} y \mathrm{~d} z+y\left(z^{2}+e^{z}\right) \mathrm{d} z \mathrm{~d} x+\left(z f(x, y)-2 e^{z}\right) \mathrm{d} x \mathrm{~d} y
\]

求 \(f(x, y)\) ，其中 \(\sum: x^{2}+y^{2}+z^{2}=1, z \geq 0\) ，方向取上侧．

\subsection{题目5}
5．设 \(a_{n} \geq 0, a_{n+1}-a_{n} \leq \int_{n}^{n+1} \frac{1}{t \ln ^{2} t+\sin ^{2} t} \mathrm{~d} t\) ，判断 \(\left\{a_{n}\right\}\) 是否收敛。

\subsection{题目6}
6．设 \(f(x)\) 在 \([a, b]\) 上非负且连续可导，证明：

\[
\left|\int_{a}^{b} f^{3}(x) \mathrm{d} x-f^{2}(a) \int_{a}^{b} f(x) \mathrm{d} x\right| \leq 2 \max _{x \in[a, b]}\left|f^{\prime}(x)\right|\left(\int_{a}^{b} f(x) \mathrm{d} x\right)^{2}
\]

\subsection{题目7}
7．设 \(f(x)\) 在 \(x=0\) 的某邻域内连续可导，且 \(\lim _{x \rightarrow 0} \frac{f(x)}{x}=2\) ，判断级数 \(\sum_{n=1}^{\infty}(-1)^{n+1} f\left(\frac{1}{\sqrt{n}}\right)\) 的敛散性。

\subsection{题目8}
8．设 \(D\) 是由 \(y=x^{2}, y=2 x^{2}, x y=1, x y=4\) 围成的区域，证明：

\[
\int_{\partial D} \frac{f(x y)}{x} \mathrm{~d} x=-\frac{\ln 2}{3} \int_{1}^{4} f^{\prime}(x) \mathrm{d} x .
\]

\subsection{题目9}
9．设函数列 \(\left\{f_{n}(x)\right\}\) 在闭区间 \([a, b]\) 上点态收敛到 \(f(x)\) ，且 \(f_{n}(x), f(x)\) 在 \([a, b]\) 上连续，\(\left\{f_{n}(x)\right\}\) 关于 \(n\) 单调，证明：\(\left\{f_{n}(x)\right\}\) 在 \([a, b]\) 上一致收敛于 \(f(x)\) ．

\subsection{题目10}
10．设 \(f(x, y)= \begin{cases}\frac{2 x y^{2}}{x^{2}+y^{4}}, & (x, y) \neq(0,0) ; \\ 0, & (x, y)=(0,0) .\end{cases}\)

\begin{itemize}
  \item[（1）] 求 \(f(x, y)\) 的偏导数．
  \item[（2）] 求 \(f(x, y)\) 在 \((0,0)\) 处的方向导数．
  \item[（3）] \(f(x, y)\) 在 \((0,0)\) 处是否可微？
\end{itemize}

\subsection{题目11}
11．解答如下问题：

\begin{itemize}
  \item[（1）] 叙述反常积分 \(\int_{a}^{+\infty} f(x, y) g(x, y) \mathrm{d} x\) 关于 \(y\) 在区间 \(I\) 上的 Abel 判别法．
  \item[（2）] \(\int_{0}^{+\infty} \frac{\cos (2 x)}{x+y} e^{-x y} \mathrm{~d} x\) 在 \(y \in[0, b]\) 上是否一致收敛？
\end{itemize}
\end{document}
